% GENERATED FILE. Edit canonical lessons, then run npm run generate:books.
% canonical-source-hash: fnv1a64:8b83aa880364ec8c
% canonical-lessons: ZH-C306-xinniang, ZH-C306-jingguo, ZH-W306-jing-pass, ZH-C306-xinren, ZH-C306-juxing, ZH-C306-xitang

\chapter{The Wedding}
\label{ch:zh-the-wedding}
\hlchaptermodality{306}

\begin{chapteropening}
\textbf{By the end of this chapter:} \emph{I can talk about the bride and the couple, the wedding sweets, and how the day went.}

Complete the last lesson of chapter 306: I can talk about the bride and the couple, the wedding sweets, and how the day went.
\end{chapteropening}

\section[xīnniáng]{xīnniáng — bride}

\label{lesson:ZH-C306-xinniang}



\begin{quote}
Before the new one: say the Mandarin for a happy occasion, then the Mandarin for the bridegroom.
\end{quote}

\subsection*{You'll want to know: xīnniáng}
\emph{xīnniáng} — \textquotedblleft{}bride\textquotedblright{}.

\begin{quote}
\emph{xīnniáng jīntiān tèbié piàoliang} — The bride is especially beautiful today.
\end{quote}

Its characters come later; for now it is a word for the ear.

\begin{cousinweb}
\emph{xīnniáng} is built from \zh{新} 'new' and \emph{niáng} 'a young woman', both inherited from Middle Chinese, \emph{sin} and \emph{nrjang} (Guangyun, in Baxter's transcription), which glosses the second as a title for a young woman.
\end{cousinweb}

\begin{grammarlens}[title={What you've built}]
One more word for telling what happened.
\end{grammarlens}

\subsection*{Your turn}
\begin{itemize}
\raggedright
  \item \emph{Say it:} \emph{xīnniáng}
  \item \emph{Say it:} \emph{xīnniáng}, once more
  \item \emph{Say it:} \emph{xīnniáng jīntiān tèbié piàoliang} — The bride is especially beautiful today.
  \item \emph{Recall:} say the Mandarin for a happy occasion, then the Mandarin for the bridegroom, then say \emph{xīnniáng} again
\end{itemize}

\begin{tcolorbox}[breakable,colback=teal!4,colframe=teal!35!black,title={Before you move on}]
What does xīnniáng mean? (\textquotedblleft{}Bride\textquotedblright{}.) Say it once more.
\end{tcolorbox}
\section[\texorpdfstring{\zh{经过} (jīngguò)}{经过 (jīngguò)}]{\zh{经过} (jīngguò) — the sequence of events; what happened}

\label{lesson:ZH-C306-jingguo}



\begin{quote}
Before the new one: say the Mandarin for the bridegroom, then the Mandarin for the bride.
\end{quote}

\subsection*{You'll want to know: \zh{经过}}
\textbf{\zh{经过}} — \emph{jīngguò} — \textquotedblleft{}the sequence of events; what happened\textquotedblright{}.

\begin{quote}
\textbf{\zh{他说了那天的经过}} \emph{tā shuō le nà tiān de jīngguò} — He told us what happened that day.
\end{quote}

The new character, in pieces: \zh{经} — eight strokes in three pieces: the silk piece of \zh{给} on the left, three strokes; on the right a fork of two strokes, a bar that turns into a sweep to the lower left and a short stroke falling to the right, over \zh{工}, three strokes. The next lesson puts it in your hand.

\begin{cousinweb}
\textbf{\zh{经过}} (\emph{jīngguò}) is built from \zh{经} 'to pass through' and \zh{过} 'to go past', both inherited from Middle Chinese, \emph{keng} and \emph{kwaH} (Guangyun, in Baxter's transcription), which glosses the first 'the warp of a cloth' and 'a path' and the second 'to cross' and 'to go beyond'.
\end{cousinweb}

\begin{grammarlens}[title={What you've built}]
One more word for telling what happened.
\end{grammarlens}

\subsection*{Your turn}
\begin{itemize}
\raggedright
  \item \emph{Say it:} \emph{jīngguò}
  \item \emph{Say it:} \emph{jīngguò}, once more
  \item \emph{Say it:} \emph{tā shuō le nà tiān de jīngguò} — He told us what happened that day.
  \item \emph{Recall:} say the Mandarin for the bridegroom, then the Mandarin for the bride, then say \emph{jīngguò} again
\end{itemize}

\begin{tcolorbox}[breakable,colback=teal!4,colframe=teal!35!black,title={Before you move on}]
What does \zh{经过} mean? (\textquotedblleft{}The sequence of events; what happened\textquotedblright{}.) Say it once more.
\end{tcolorbox}
\section[\texorpdfstring{\zh{经} (jīng)}{经 (jīng)}]{\zh{经} — the silk piece, then a fork over \zh{工}}

\label{lesson:ZH-W306-jing-pass}



\begin{quote}
The word first, then the character inside it.

\begin{itemize}
\raggedright
  \item \emph{Recall:} say \emph{jīngguò} — \textbf{R1}, one lesson back, before it decays
\end{itemize}

\begin{quote}
\zh{给} \zh{工}
\end{quote}

\emph{Read it:} each one aloud, with nothing said first — \emph{gěi}, \emph{gōng}
\end{quote}

\subsection*{Script — the silk piece, then a fork over \zh{工}}
\hlblockfigure{\detokenize{figures/ZH-W306-jing-pass-filmstrip.pdf}}{How \zh{经} is written, stroke by stroke}

\begin{quote}
\zh{经}
\end{quote}

\textbf{Three pieces: the silk piece on the left; on the right a fork over \zh{工}.} Count the pieces, not the strokes.

\textbf{The silk piece, three strokes.} As on the left of \zh{给}: two strokes that each drop to the left and turn to the right, then a stroke rising to the right.

\textbf{The fork, two strokes.} A bar that turns into a sweep down to the lower left, then a short stroke falling to the right.

\textbf{\zh{工}, three strokes}: a short bar, a vertical, and a long bar across the foot.

3 + 2 + 3 = eight.

\subsection*{Your turn}
\begin{itemize}
\raggedright
  \item \emph{Write it:} \zh{经}, all eight, with the model in front of you
  \item \emph{Cover:} write it again from memory, and count the pieces, not the strokes
\end{itemize}

\begin{tcolorbox}[breakable,colback=teal!4,colframe=teal!35!black,title={Before you move on}]
How many pieces in \zh{经}? (\textbf{Three}.) Where have you made the left piece? (\textbf{On the left of \zh{给}}.) What sits at the foot of the right side? (\textbf{\zh{工}}.)

Source: \href{https://raw.githubusercontent.com/chanind/hanzi-writer-data/68d10a4b21150cae5e1ebbd223eed289cf32d90c/data/\%E7\%BB\%8F.json}{Hanzi Writer Data \zh{经}, pinned stroke paths and medians}.
\end{tcolorbox}
\section[\texorpdfstring{\zh{新人} (xīnrén)}{新人 (xīnrén)}]{\zh{新人} (xīnrén) — the bride and groom}

\label{lesson:ZH-C306-xinren}



\begin{quote}
Before the new one: say the Mandarin for the bride, then the Mandarin for the whole story of what happened.
\end{quote}

\subsection*{You'll want to know: \zh{新人}}
\textbf{\zh{新人}} — \emph{xīnrén} — \textquotedblleft{}the bride and groom\textquotedblright{}.

\begin{quote}
\textbf{\zh{今天的新人是我的同学}} \emph{jīntiān de xīnrén shì wǒ de tóngxué} — The couple getting married today were my classmates.
\end{quote}

\begin{cousinweb}
\textbf{\zh{新人}} (\emph{xīnrén}) is built from \zh{新} 'new' and \zh{人} 'a person', both inherited from Middle Chinese, \emph{sin} and \emph{nyin} (Guangyun, in Baxter's transcription): the new people, the couple who marry that day.
\end{cousinweb}

\begin{grammarlens}[title={What you've built}]
One more word for telling what happened.
\end{grammarlens}

\subsection*{Your turn}
\begin{itemize}
\raggedright
  \item \emph{Say it:} \emph{xīnrén}
  \item \emph{Say it:} \emph{xīnrén}, once more
  \item \emph{Say it:} \emph{jīntiān de xīnrén shì wǒ de tóngxué} — The couple getting married today were my classmates.
  \item \emph{Recall:} say the Mandarin for the bride, then the Mandarin for the whole story of what happened, then say \emph{xīnrén} again
\end{itemize}

\begin{tcolorbox}[breakable,colback=teal!4,colframe=teal!35!black,title={Before you move on}]
What does \zh{新人} mean? (\textquotedblleft{}The bride and groom\textquotedblright{}.) Say it once more.
\end{tcolorbox}
\section[jǔxíng]{jǔxíng — to hold a ceremony; to hold an event}

\label{lesson:ZH-C306-juxing}



\begin{quote}
Before the new one: say the Mandarin for the whole story of what happened, then the Mandarin for the bride and groom.
\end{quote}

\subsection*{You'll want to know: jǔxíng}
\emph{jǔxíng} — \textquotedblleft{}to hold a ceremony; to hold an event\textquotedblright{}.

\begin{quote}
\emph{hūnlǐ zài zhōngwǔ jǔxíng} — The wedding is at noon.
\end{quote}

Its characters come later; for now it is a word for the ear.

\begin{cousinweb}
\emph{jǔxíng} is built from \emph{jǔ} 'to raise' and \emph{xíng} 'to go', both inherited from Middle Chinese, \emph{kjoX} and \emph{haeng} (Guangyun, in Baxter's transcription), which glosses the first 'to lift' and 'to set moving' and the second 'to walk': to set going.
\end{cousinweb}

\begin{grammarlens}[title={What you've built}]
One more word for telling what happened.
\end{grammarlens}

\subsection*{Your turn}
\begin{itemize}
\raggedright
  \item \emph{Say it:} \emph{jǔxíng}
  \item \emph{Say it:} \emph{jǔxíng}, once more
  \item \emph{Say it:} \emph{hūnlǐ zài zhōngwǔ jǔxíng} — The wedding is at noon.
  \item \emph{Recall:} say the Mandarin for the whole story of what happened, then the Mandarin for the bride and groom, then say \emph{jǔxíng} again
\end{itemize}

\begin{tcolorbox}[breakable,colback=teal!4,colframe=teal!35!black,title={Before you move on}]
What does jǔxíng mean? (\textquotedblleft{}To hold a ceremony; to hold an event\textquotedblright{}.) Say it once more.
\end{tcolorbox}
\section[xǐtáng]{xǐtáng — wedding sweets}

\label{lesson:ZH-C306-xitang}



\begin{quote}
Before the new one: say the Mandarin for the bride and groom, then the Mandarin for to hold a ceremony.
\end{quote}

\subsection*{You'll want to know: xǐtáng}
\emph{xǐtáng} — \textquotedblleft{}wedding sweets\textquotedblright{}.

\begin{quote}
\emph{wǒmen chī le tāmen de xǐtáng} — We had sweets at their wedding.
\end{quote}

Its characters come later; for now it is a word for the ear.

\begin{cousinweb}
\emph{xǐtáng} is built from \zh{喜} 'joy' and \emph{táng} 'sugar', both inherited from Middle Chinese, \emph{xiX} and \emph{dang} (Guangyun, in Baxter's transcription), which glosses the second 'a sweet syrup'.
\end{cousinweb}

\begin{grammarlens}[title={What you've built}]
One more word for telling what happened.
\end{grammarlens}

\subsection*{Your turn}
\begin{itemize}
\raggedright
  \item \emph{Say it:} \emph{xǐtáng}
  \item \emph{Say it:} \emph{xǐtáng}, once more
  \item \emph{Say it:} \emph{wǒmen chī le tāmen de xǐtáng} — We had sweets at their wedding.
  \item \emph{Recall:} say the Mandarin for the bride and groom, then the Mandarin for to hold a ceremony, then say \emph{xǐtáng} again
\end{itemize}

\begin{tcolorbox}[breakable,colback=teal!4,colframe=teal!35!black,title={Before you move on}]
What does xǐtáng mean? (\textquotedblleft{}Wedding sweets\textquotedblright{}.) Say it once more.
\end{tcolorbox}
